\section{Introduction}
\label{sec:intro}

\Cref{lst:motivating} shows a loop nest from the dynamical core of the ICON
climate model~\cite{icon}, lightly distilled. Three kernels---\texttt{A},
\texttt{B}, and \texttt{C}---sweep the same column of edge values, reusing the
horizontal and tangential velocities \texttt{vn} and \texttt{vt}. A production
compiler keeps them as three separate loop nests: it cannot fuse \texttt{A} with
\texttt{C} across the opaque \lstinline[style=fortran]|IF (.NOT. lvn_only)| guard,
and it cannot drop the \lstinline[style=fortran]|IF (istep == 1)| test, because
both \texttt{lvn\_only} and \texttt{istep} are read from the model's configuration
file---its \emph{namelist}---at runtime. So on a memory-bandwidth-bound code, the
velocities are streamed from memory three times where once would do.

\begin{lstlisting}[style=fortran,float=t,caption={A configuration-dependent fusion
barrier in ICON's \texttt{velocity\_tendencies} (distilled). \texttt{istep} and
\texttt{lvn\_only} are read once from the namelist at startup and never change,
yet a compiler must treat them as unknown.},label={lst:motivating}]
! istep, lvn_only: read once from the namelist, then never change.
IF (istep == 1) THEN
  DO jk = 2, nlev                          ! kernel A
    vn_ie(je,jk)       = wgtfac_e*vn + (1-wgtfac_e)*vn_km1
    z_kin_hor_e(je,jk) = 0.5*(vn**2 + vt**2)
  END DO
  IF (.NOT. lvn_only) THEN                 ! <-- config-guarded fusion barrier
    DO jk = 2, nlev                        ! kernel B
      z_vt_ie(je,jk)   = wgtfac_e*vt + (1-wgtfac_e)*vt_km1
    END DO
  END IF
  DO jk = 2, nlev                          ! kernel C
    z_w_concorr_me(je,jk) = vn*ddxn_z + vt*ddxt_z
  END DO
END IF
\end{lstlisting}

These values are constants of the run. A weather center launches this code with a
fixed namelist and integrates for hours or days without ever changing
\texttt{lvn\_only} or \texttt{istep}. The moment they are known the barriers
dissolve: with \texttt{istep}\,$=$\,2 the block is dead; with
\texttt{istep}\,$=$\,1 and \texttt{lvn\_only}\,$=$\,\texttt{.true.}, kernel
\texttt{B} vanishes and \texttt{A} and \texttt{C} fuse into one column sweep; with
\texttt{lvn\_only}\,$=$\,\texttt{.false.}, \texttt{B} becomes unconditional and all
three fuse. The information that unlocks every one of these optimizations is
sitting, in plain text, in a file on disk.

\paragraph{Why the compiler cannot reach it.}
One might expect aggressive optimization to recover this. It cannot, for two
structural reasons rather than a missing pass, which we make precise in
\Cref{sec:background} on a NAS~LU kernel whose grid bounds are module variables (as in
any code built from \texttt{COMMON} blocks or configuration modules). First, the value
is \emph{not in the program at all}: it enters through an \emph{I/O barrier}, a
namelist read that lowers to opaque runtime calls past which constant propagation must
assume anything. Second, it is \emph{module-level state}---external by linkage, reached
across hundreds of translation units through
\lstinline[style=fortran]|USE ... ONLY|---so a single translation unit cannot even
assume it is unwritten. Production configuration state fails on both counts, so it
stays a runtime value at every optimization level, including whole-program LTO
(\Cref{tab:lu}). The cost is a \emph{configuration tax}: branches, dead code, and
fusion and vectorization inhibitors that survive every optimization level.

\paragraph{The configuration file is the profile.}
Profile-guided optimization specializes for values that are stable in practice by
instrumenting a run. We observe that for configuration variables no profiling is
needed: \emph{the configuration file already is the profile}, and an exact one.
This reframes specialization around a third point in the compilation lifecycle.
Ahead-of-time compilation does not know the deployment values; just-in-time
compilation knows them but pays a runtime cost on every execution;
\emph{deployment-time specialization} knows the values---from the config file---and
pays no runtime cost at all. We parse the namelist, trace each configuration
variable across the I/O barrier to the glacial module state it initializes, recover
its domain, and inject it. For the scalar case this is precisely the missing input:
once the value is supplied \emph{as a source-level constant}, the standard cascade
does the rest (the regime \Cref{tab:lu} probes). The contribution there is not the
cascade but the bridge---recovering, statically and soundly, a value that lives in a
data file the compiler never reads.

\td{DEEP ARG B (PLAN.md 2.B) --- the config file is an EXACT, GUARD-FREE profile.
Value-knowledge spectrum: source-const(parse) | other-TU(LTO) | config/glacial
(deployment, EXACT) | PGO(post-run, probabilistic, NEEDS GUARD) | runtime. RIFS must
guard because values are merely frequent; we never guard because the namelist is
authoritative. Chain: exact $\Rightarrow$ no per-call guard $\Rightarrow$ unguarded
whole-program specialization $\Rightarrow$ covering set is the natural deployment unit
(ties C1/C2/C3). Rename the hook to "exact, guard-free profile".}
\td{DEEP ARG C (PLAN.md 2.C) --- automatic staging (dissolves "just partial
evaluation"). A configured code is a TWO-STAGE program; PE/SFAC need a hand-supplied
binding-time annotation, we DISCOVER the boundary automatically and EXECUTE stage 1 at
deployment-compile-time. CAVEAT (skeptical review): the boundary is NOT the textual
READ (glacial fixpoint reached later via config derivation / MPI broadcast / restart
override) --- define it as the last point dominating timestep-loop entries with empty
compute-phase may-write set (detail in \Cref{sec:recovery}). Elevates Autrey--Wolfe +
extends to types.}

\paragraph{Beyond scalars: interval arithmetic on collections.}
A recovered value, however, is rarely just a scalar, and the cases a standard
pipeline cannot finish even with the value in hand are the ones that matter most.
We treat four domain kinds uniformly under one lattice we call \emph{interval
arithmetic on collections}: a \emph{singleton} folds to a constant; a small
\emph{finite value-set} (a boolean, a few integer modes) drives specialization over
the cross-product of its elements; an \emph{integer interval}, recovered from the
shape an index array is used to address, narrows that array to the smallest integer
type that fits; and a finite \emph{type-set} collapses dynamic dispatch. The last is
our central mechanism. A configuration variable that selects one of several linear
solvers has a finite type-set; lifting the recovered value to that type-set replaces
the v-table with an \lstinline[style=fortran]|if|-chain that configuration constants
prune to the single deployed solver. Neither index narrowing nor devirtualization
follows from injecting a scalar and re-running the optimizer; indeed the default
Fortran pipeline does not devirtualize from a known dynamic type at all
(\Cref{sec:spec}). To our knowledge no prior specializer derives a \emph{type}-set
from a recovered configuration \emph{value} to drive devirtualization
(\Cref{sec:related}).

\paragraph{Why this matters.}
Recovering these gains by hand is a recognized, expensive craft. ECMWF's recent
twelvefold acceleration of the ecRad radiation scheme~\cite{ecrad2024} reports a
separable threefold ``code optimization'' factor; hard-coding the spectral g-point
count into build flags so that loop bounds become constant---manual configuration
propagation---is one of the hand optimizations that contribute to it. Such work
costs hundreds of person-hours, must be redone per configuration, and is fragile: a
mismatched setting silently corrupts or crashes the model. Deployment-time
specialization automates it, and because it knows the dimensions at compile time it
also lets the GPU backend emit tighter kernels---constant launch bounds, unrolled
loops, narrowed index types---that a runtime-valued parameter forecloses.

\td{DEEP ARG E (PLAN.md 2.E) --- systems value: production codes are generic for
MAINTAINABILITY/DISTRIBUTION, not because any single deployment needs generality
(every deployment runs a fixed config), so the genericity is pure runtime overhead
(the configuration tax). The real proposition is recovering specialized performance
WITHOUT sacrificing generic-source maintainability --- what the manual ecRad route
cannot do. Deeper than any single speedup number; state it here.}

\paragraph{Contributions.}
\begin{itemize}
  \item We characterize the \emph{configuration tax} and show, on a NAS~LU kernel,
  that the barrier is structural for two independent reasons---value-behind-I/O and
  external linkage---so no optimization level recovers production configuration state
  (\Cref{sec:background}). We position \emph{deployment-time specialization}
  (config-file-driven AOT specialization with startup binary selection) against AOT,
  JIT, and PGO.

  \item \textbf{Glacial recovery across the I/O barrier} (\Cref{sec:recovery}): a
  static analysis that traces namelist reads to the glacial module-level state they
  initialize and recovers it as deployment-time constants---the value the standard
  pipeline is missing and structurally cannot obtain.

  \item \textbf{Type-domain devirtualization} (\Cref{sec:spec}, our centerpiece):
  we lift the value/interval lattice to a finite type-set; to our knowledge the
  first analysis to drive devirtualization from a configuration-recovered value
  domain, eliminating dispatch that the default Fortran pipeline (flang, nvfortran)
  does not remove even when the dynamic type is known.

  \item \textbf{The variant economy} (\Cref{sec:variants}): finite-domain
  cross-product specialization with greedy minimal-covering-set selection over a
  distribution of operational namelists, deployed as multiple \emph{unguarded}
  whole-program binaries dispatched once at startup.

  \item \textbf{Evaluation} (\Cref{sec:eval}) on ICON, ecRad, and Quantum ESPRESSO
  across GH200 and MI300A, against \texttt{nvfortran -O4}\footnote{Current nvfortran
  (26.3) \emph{deprecates and ignores} \texttt{-Mipa}, so even \texttt{-O4} performs
  no interprocedural constant propagation across the namelist read.} and against the
  strongest standard-tool baseline---the recovered values injected as constants
  under whole-program LTO, which captures the scalar cases but not the
  devirtualization and index narrowing.
\end{itemize}

We build on the DaCe dataflow compiler and its Fortran frontend, which gives the
whole-program view our specializations need without depending on the maturity of
Flang/LLVM LTO on production-scale Fortran.
